\section{結論與研究限制}

\subsection{研究結論}

本研究將 YouTube 留言視為社群互動資料，提出並實作一套以「留言社群洞察」為核心的 audience intelligence
分析框架，整合留言者涵蓋率分層、回訪率、co-commenter network 社群偵測、Qwen 情緒與 ABSA、reply-thread
conflict 與外部事件時間軸，並以 48 個台灣頻道作為 broad benchmark。以 The DoDo Men 為個案，研究發現可
歸納如下：

\begin{itemize}
  \item \textbf{（RQ1）} 頻道留言社群整體健康，屬「穩定討論引擎型」：跨影片回訪率（PR~81）與近期回訪率
        （PR~91）高於同儕、一次性留言者占比低於同儕（PR~15），互動主要來自穩定核心與回訪觀眾，而非一次性流量。
  \item \textbf{（RQ2）} 頻道存在三個觀眾社群——知識職涯型（47.7\%，最穩定、負面最低）、挑戰冒險型
        （32.0\%，高投入但負面最高）與旅遊來賓型（20.3\%）；惟 modularity 僅 0.31，社群為軟性偏好傾向而非
        壁壘分明的派系，且高度集中（社群集中度指數 87）。
  \item \textbf{（RQ3）} 正面情緒主要針對主持人與內容品質，負面情緒主要針對價值觀／政治與真實性，屬較高層級
        的聲譽風險；少數影片的負面被按讚顯著放大（如某影片按讚加權負面率達 49.3\%）。
  \item \textbf{（RQ4）} 負面與衝突並非同一件事：多數衝突為「對立母串」（回覆與母留言立場相反）而非「圍剿」，
        少數影片（如建中教課）同時具高圍剿與高對立母串，屬需優先關注的留言串。
  \item \textbf{（RQ5）} 僅少數外部事件與 YouTube 留言變化出現時間上的關聯跡象，整體而言外部討論並未顯著、
        穩定地帶入可長期經營的新留言者。
\end{itemize}

整體而言，本研究示範了如何從「互動數量」進一步分析「互動品質」，將高留言、高互動拆解為更具策略意義的問題，
並轉化為內容策略、分眾溝通策略與炎上監測策略。

\subsection{研究限制}
\label{sec:limit}

\begin{enumerate}
  \item 本研究分析對象為 active commenting audience，而非全體觀看者。許多觀眾只觀看不留言，因此留言社群
        不等同於整體觀眾結構。
  \item co-commenting relationship 不代表真實社交關係。兩位留言者留言同一支影片，只能說明他們共同參與相同
        內容，不代表兩者彼此認識或互相影響。
  \item YouTube API 與資料抓取限制可能導致留言資料不完整，例如刪除留言、不可見留言或歷史回溯限制。
  \item Qwen 主題與情緒標註可能受到反諷、迷因、口語用法與上下文缺失影響，因此結果需搭配人工檢查或抽樣驗證。
  \item reply-thread conflict 指標只能反映留言串中的情緒對立與互動結構，不能完全等同於真實衝突或社群傷害。
  \item 外部事件分析只能解讀為關聯跡象，而非因果推論。事件後 YouTube 留言變化可能同時受到影片發布、演算法
        推薦、媒體報導或其他未觀察因素影響；本研究外部資料（PTT／Dcard）亦相對稀疏。
  \item 本研究未使用 YouTube Studio 後台資料，因此無法分析觀看序列、觀眾留存曲線、推薦來源或實際轉換行為。
  \item 本研究 benchmark 使用 48 個台灣 creator／media 頻道作為 broad benchmark，尚未依頻道主題、受眾
        組成或內容格式建立 matched peer group，因此百分位數僅作為相對位置與異常程度參考，不應被解讀為同類型
        頻道排名。
  \item ABSA 分析目前未設置對照組（control group），因此 aspect 層級的負面比較僅為頻道內部相對結構，尚難
        與外部基準對照。
\end{enumerate}

\subsection{未來研究}

未來研究可朝四個方向延伸。第一，建立 \textbf{matched peer cohort}：依頻道主題（如旅遊／雙主持類）、受眾組成
與內容格式挑選對照頻道，使 benchmark 百分位更具可比性。第二，為 ABSA 與情緒分析設置\textbf{對照組與抽樣人工
驗證}，以評估標註可靠度並支援 aspect 層級的跨頻道比較。第三，將外部事件分析擴展至更多來源（如 Threads、新聞）
並嘗試更嚴謹的\textbf{準實驗設計}（如 difference-in-differences），以更穩健地檢驗外部聲量與 YouTube 留言
變化的關聯。第四，將本框架\textbf{產品化}為可重複套用的 dashboard，使創作者能持續監測留言社群健康、內容敏感度
與炎上風險。
